% P15-U1 / N-x-008: integration and expectation foundations.
\chapter{Integration and Expectation}
\label{sec:p15-integration-expectation}

An integral extends a weighted finite sum to measurable functions. The
extension must specify when an infinite value is allowed, when a signed
average is undefined, and which limits preserve an average. This chapter
constructs that integral and proves the convergence and expectation results
needed for subsequent probability theory.

The source is Sections 1.4--1.6 of Durrett's author version of
\emph{Probability: Theory and Examples}, dated 11 January 2019
\cite[pp.~18--37]{DurrettPTEAuthor2019}. We begin with nonnegative simple
functions and their supremum, prove monotone convergence directly, and
then pass to signed functions. A proved comparison below identifies this
integral with Durrett's four-stage construction on sigma-finite spaces.
The event, measure, approximation and real-line constructions used here
are in Chapters~\ref{sec:p15-measurable-probability}
and~\ref{sec:p15-measure-extension}. No integral is used to construct
the measure with respect to which it is taken.

\section{Nonnegative integration}
\label{sec:p15-nonnegative-integral}

\begin{hbdefinition}{def:p15-simple-integral}{The integral of a nonnegative simple function}
Fix a measure space \((\Omega,\mathcal F,\mu)\). Unless a theorem
says otherwise, neither finiteness nor sigma-finiteness is assumed.
A nonnegative simple function has a representation
\(s=\sum_{j=1}^m a_j\mathbf1_{A_j}\), with \(m\geq1\), where the
\(A_j\) are pairwise disjoint measurable sets with union \(\Omega\),
and \(a_j\) are finite
nonnegative real numbers. Empty parts and repeated values are harmless.
An omitted zero region can be added as a part with coefficient zero.
Define
\[
 \int s\,d\mu\defeq\sum_{j=1}^m a_j\mu(A_j).
\]
In nonnegative extended arithmetic we use \(0\cdot\infty=\infty\cdot0=0\).
For a positive finite \(a\), \(a\cdot\infty=\infty\).
Also \(\infty\cdot\infty=\infty\).
This makes the contribution of a zero-valued region zero even when that
region has infinite measure. The next theorem proves independence of
the chosen representation.
\index{integral!simple function}\index{integration!zero times infinity}
\end{hbdefinition}

\begin{hbtheorem}{thm:p15-simple-integral-laws}{proved=annex:p15.simple-integral-laws}{Simple integrals and refinement}
The simple integral is well defined. For nonnegative simple \(s,t\)
and finite \(c\geq0\),
\[
 \int(s+t)\,d\mu=\int s\,d\mu+\int t\,d\mu,\qquad
 \int cs\,d\mu=c\int s\,d\mu.
\]
If \(s\leq t\), then \(\int s\,d\mu\leq\int t\,d\mu\).
In particular \(\int\mathbf1_A\,d\mu=\mu(A)\). A finite sum
\(s=\sum_j a_j\mathbf1_{B_j}\) with nonnegative coefficients has
integral \(\sum_j a_j\mu(B_j)\) even when the \(B_j\) overlap.
\end{hbtheorem}

\begin{hbdefinition}{def:p15-nonnegative-integral}{The nonnegative integral and almost-everywhere statements}
For measurable \(f:\Omega\to[0,\infty]\), define
\[
 \int f\,d\mu\defeq
 \sup\left\{\int s\,d\mu:s\text{ nonnegative simple},\ 0\leq s\leq f\right\}.
\]
The zero function is an admissible lower approximation, so the set is
nonempty. This integral always exists in \([0,\infty]\).
For \(A\in\mathcal F\), write
\(\int_A f\,d\mu\defeq\int f\mathbf1_A\,d\mu\), where
\(f\mathbf1_A\) equals \(f\) on \(A\) and zero outside \(A\),
including where \(f=\infty\).

A statement holds \(\mu\)-\emph{almost everywhere}, abbreviated
\emph{a.e.}, if it holds outside a measurable set of measure zero.
Every function integrated in this chapter must be measurable. Statements
about replacing a function on a null set refer to replacements that are
also measurable; completeness is not silently assumed.
\index{integral!nonnegative}\index{almost everywhere}
\end{hbdefinition}

\begin{hbtheorem}{thm:p15-monotone-integral}{proved=annex:p15.monotone-integral}{Monotone convergence}
The nonnegative integral agrees with the simple integral and is monotone:
\(0\leq f\leq g\) implies \(\int f\,d\mu\leq\int g\,d\mu\).
If nonnegative measurable functions satisfy \(f_n\leq f_{n+1}\)
everywhere and \(f_n\to f\) pointwise, then
\[
 \int f_n\,d\mu\ \uparrow\ \int f\,d\mu.
\]
The limit and the integral may be infinite.
\end{hbtheorem}

\begin{hbtheorem}{thm:p15-nonnegative-integral-laws}{proved=annex:p15.nonnegative-integral-laws}{Nonnegative calculus, null sets and Fatou's lemma}
\par\nobreak
For nonnegative measurable \(f,g,f_n\) and finite \(c\geq0\),
\[
 \begin{aligned}
 \int(f+g)\,d\mu&=\int f\,d\mu+\int g\,d\mu,\\
 \int cf\,d\mu&=c\int f\,d\mu,\\
 \int\sum_{n\geq1}f_n\,d\mu&=\sum_{n\geq1}\int f_n\,d\mu.
 \end{aligned}
\]
All sums are nonnegative extended sums. Moreover:
\begin{enumerate}
\item \(\int f\,d\mu=0\) exactly when \(f=0\) a.e.; if
      \(\int f\,d\mu<\infty\), then \(f<\infty\) a.e.
\item If \(f\leq g\) a.e., their integrals have the same order.
      If \(f=g\) a.e., their integrals agree. Monotone convergence
      therefore also holds when its pointwise hypotheses hold a.e.
\item Fatou's inequality is
      \(\int\liminf_n f_n\,d\mu\leq\liminf_n\int f_n\,d\mu\).
\item The function \(\nu_f(A)\defeq\int_A f\,d\mu\), for
      \(A\in\mathcal F\), is a measure. In particular an integral
      over a disjoint measurable union is the sum of the integrals
      over its members.
\end{enumerate}
\end{hbtheorem}

\section{Signed integration and finite approximations}
\label{sec:p15-signed-integral}

\begin{hbdefinition}{def:p15-integrable-functions}{Positive and negative parts, integrability and the signed integral}
Negation on the extended reals exchanges the two infinite endpoints.
For a measurable extended real function \(f\), define pointwise
\[
 f^+\defeq\max(f,0),\qquad f^-\defeq\max(-f,0),\qquad
 |f|\defeq f^++f^-.
\]
These are nonnegative measurable functions. Its signed integral is
\[
 \int f\,d\mu\defeq\int f^+\,d\mu-\int f^-\,d\mu
\]
provided at least one of the two nonnegative integrals is finite.
It may then equal \(\infty\) or \(-\infty\). If both are infinite,
the signed integral is undefined; it is not assigned a value by cancellation.
The function is \emph{integrable} when \(\int|f|\,d\mu<\infty\).
Then both part integrals and the signed integral are finite.
For integrable \(f\) and \(A\in\mathcal F\), the notation
\(\int_A f\,d\mu\) uses \(f\) on \(A\) and zero on \(A^c\).
This restricted function is integrable because its absolute value is
at most \(|f|\).

An integrable extended real function is finite a.e. by the preceding
theorem. Replacing its values on the measurable null set
\(\{|f|=\infty\}\) by zero gives a real-valued measurable
representative with the same integral. Linearity statements for integrable
functions below use real-valued representatives, so no pointwise sum
\(\infty+(-\infty)\) is formed.
\index{integrable function}\index{integral!signed}\index{positive and negative parts}
\end{hbdefinition}

\begin{hbtheorem}{thm:p15-signed-integral-laws}{proved=annex:p15.signed-integral-laws}{Linearity and the absolute-value bound}
\par\nobreak
If real measurable \(f,g\) are integrable and \(a,b\in\mathbb R\),
then \(af+bg\) is integrable and
\[
 \int(af+bg)\,d\mu=a\int f\,d\mu+b\int g\,d\mu,
 \qquad
 \left|\int f\,d\mu\right|\leq\int|f|\,d\mu.
\]
If \(f\leq g\) a.e., then \(\int f\,d\mu\leq\int g\,d\mu\).
Measurable a.e. replacement preserves integrability and the integral.
If \(f=u-v\) with nonnegative real integrable \(u,v\), then
\(\int f\,d\mu=\int u\,d\mu-\int v\,d\mu\).
For an integrable \(f\) and disjoint measurable sets \(A_n\),
\[
 \int_{\bigcup_n A_n}f\,d\mu=\sum_n\int_{A_n}f\,d\mu,
 \qquad
 \sum_n\left|\int_{A_n}f\,d\mu\right|<\infty.
\]
\end{hbtheorem}

\begin{hbtheorem}{thm:p15-integral-source-bridge}{proved=annex:p15.integral-source-bridge}{Finite-measure approximation and the source construction}
\par\nobreak
Let \(E\in\mathcal F\) have finite measure, and let real measurable
\(f\) vanish outside \(E\) and satisfy \(|f|\leq M\), where \(0\leq M<\infty\).
Then \(f\) is integrable and
\[
 \int f\,d\mu
 =\sup_{s\leq f}\int s\,d\mu
 =\inf_{t\geq f}\int t\,d\mu,
\]
where \(s,t\) range over real simple functions vanishing outside \(E\).
Separately, let \(f:\Omega\to[0,\infty]\) be any nonnegative
measurable function. If \(E_n\uparrow\Omega\) with
\(\mu(E_n)<\infty\), then
\[
 \int_{E_n}\min(f,n)\,d\mu\ \uparrow\ \int f\,d\mu.
\]
Consequently this last integral also equals the supremum of
\(\int h\,d\mu\) over bounded nonnegative measurable \(h\leq f\)
whose positive region has finite measure. These identities identify the
present construction with Durrett's bounded and nonnegative stages on
sigma-finite spaces; their signed stages then agree by positive and
negative parts.
\end{hbtheorem}

\begin{hbdefinition}{def:p15-lp-convergence}{Integral sizes and modes of convergence}
For a real measurable function and \(1\leq p<\infty\), write
\[
 \|f\|_p\defeq\left(\int|f|^p\,d\mu\right)^{1/p},\qquad
 \|f\|_\infty\defeq
 \inf\{M\geq0:\mu(\{|f|>M\})=0\},
\]
with infimum of an empty set equal to infinity. The space \(L^1(\mu)\)
consists of integrable real measurable functions identified when they
agree a.e. Its norm is \(\|f\|_1\); the triangle inequality and
\(\|f\|_1=0\) exactly for the zero a.e. class follow from the
preceding integral laws. No completeness theorem for this space is
asserted here.

Convergence in \(L^1\) means \(\|f_n-f\|_1\to0\).
Convergence \emph{in measure} means
\(\mu(\{|f_n-f|>\varepsilon\})\to0\) for every \(\varepsilon>0\).
These definitions concern the specified measure. On a probability space,
convergence in measure is called convergence in probability.
\index{convergence!in measure}\index{convergence!in L1}\index{Lp norm}
\end{hbdefinition}

\begin{hbtheorem}{thm:p15-dominated-bounded-convergence}{proved=annex:p15.dominated-bounded-convergence}{Dominated and bounded convergence}
\par\nobreak
Let \(f_n,f\) be real measurable functions.
\begin{enumerate}
\item If \(f_n\to f\) a.e. and \(|f_n|\leq g\) a.e. for every
      \(n\), where \(g\geq0\) is integrable, then all \(f_n\)
      and \(f\) are integrable and
      \(\|f_n-f\|_1\to0\). Hence their integrals converge to
      \(\int f\,d\mu\).
\item If \(\mu(E)<\infty\), each \(f_n\) vanishes outside \(E\)
      a.e., \(|f_n|\leq M\) a.e. for \(0\leq M<\infty\), and \(f_n\to f\)
      in measure, then \(f=0\) outside \(E\) a.e., \(|f|\leq M\)
      a.e., and again \(\|f_n-f\|_1\to0\) and the integrals converge.
\end{enumerate}
In particular uniformly bounded a.e. convergence on a finite measure
space satisfies the first assertion with a constant dominating function.
\end{hbtheorem}

\begin{hbexample}{def:p15-convergence-failures}{Concentration, an infinite domain and a missing lower bound}
On the constructed uniform space \((0,1)\), put
\(f_n(u)=n\mathbf1_{(0,1/n]}(u)\). Each integral is one, while
\(f_n\to0\) everywhere. Thus Fatou's inequality is strict and
convergence of integrals fails without a suitable bound. Any measurable
nonnegative dominator \(g\) would satisfy \(g\geq k\) a.e. on
\((1/(k+1),1/k]\). These disjoint intervals therefore give
\[
 \int g\,d\mathbb U\geq
 \sum_{k\geq1}k\left(\frac1k-\frac1{k+1}\right)
 =\sum_{k\geq1}\frac1{k+1}=\infty.
\]
The last series diverges because each sufficiently large block between
successive powers of two has sum bounded below by a fixed positive
number. Hence no integrable dominator exists. Replacing \(f_n\) by
\(-f_n\) also disproves an unrestricted signed version of Fatou:
the integrals have lower limit \(-1\), whereas the limit function has
integral zero.

On \(\mathbb R\) with the constructed Lebesgue measure, the functions
\(h_n=n^{-1}\mathbf1_{(0,n]}\) are bounded by one and converge uniformly
to zero, hence in measure, but have integrals one. This removes the
finite-domain hypothesis from bounded convergence. On the same space,
\(r_n=-\mathbf1_{[n,\infty)}\) increases pointwise to zero, but each
signed integral is \(-\infty\). Monotonicity alone does not replace
the nonnegativity condition in the stated monotone convergence theorem.
\end{hbexample}

\section{Inequalities for integrals}
\label{sec:p15-integral-inequalities}

\begin{hbtheorem}{thm:p15-holder-integral}{proved=annex:p15.holder-integral}{H\"older and Cauchy--Schwarz inequalities}
\par\nobreak
For real measurable \(f,g\) and \(p,q\in[1,\infty]\) with
\(1/p+1/q=1\), using \(1/\infty=0\),
\[
 \int|fg|\,d\mu\leq\|f\|_p\|g\|_q.
\]
The product on the right uses \(0\cdot\infty=0\); all zero-norm
and infinite-norm cases are included. In particular, when both displayed
norms are finite, the product is integrable. The case \(p=q=2\) is
the Cauchy--Schwarz inequality.
\end{hbtheorem}

\begin{hbdefinition}{def:p15-convex-integral}{Convex functions}
A function \(\varphi:\mathbb R\to\mathbb R\) is \emph{convex}
when, for real \(x,y\) and \(0<t<1\),
\[
 \varphi(tx+(1-t)y)\leq t\varphi(x)+(1-t)\varphi(y).
\]
It is finite valued throughout its domain. The proof below establishes
its continuity and its supporting affine function at every real point.
\index{convex function}
\end{hbdefinition}

\begin{hbtheorem}{thm:p15-jensen-integral}{proved=annex:p15.jensen-integral}{Jensen's inequality}
Let \(\mu(\Omega)=1\), let \(f\) be real and integrable, and let
\(\varphi:\mathbb R\to\mathbb R\) be convex. Then \(\varphi(f)\)
is measurable. If it is integrable, then
\[
 \varphi\left(\int f\,d\mu\right)\leq\int\varphi(f)\,d\mu.
\]
The probability normalization and both integrability hypotheses are part
of this statement.
\end{hbtheorem}

\section{Expectation and computation from a law}
\label{sec:p15-general-expectation}

\begin{hbdefinition}{def:p15-expectation-integral}{Expectation and moments}
On a probability space \((\Omega,\mathcal F,\mathbb P)\), define
\(\mathbb E[X]\defeq\int X\,d\mathbb P\) with the preceding
existence conventions. For nonnegative \(X\), it always exists in
\([0,\infty]\). For signed \(X\), it exists exactly when the two
part expectations are not both infinite. A \emph{finite mean} requires
\(\mathbb E[|X|]<\infty\). Write
\(\mathbb E[X;A]\defeq\int_A X\,d\mathbb P\) when \(X\geq0\)
or \(X\) is integrable and \(A\in\mathcal F\).
For a positive integer \(k\), the \(k\)-th moment is
\(\mathbb E[X^k]\) when this expectation is defined.
For real \(X\) with \(\mathbb E[X^2]<\infty\), its finite mean
\(m=\mathbb E[X]\) and \emph{variance}
\(\operatorname{Var}(X)\defeq\mathbb E[(X-m)^2]\) are justified
by the results below.
\index{expectation!as an integral}\index{moment}\index{variance}
\end{hbdefinition}

\begin{hbtheorem}{thm:p15-expectation-law-integral}{proved=annex:p15.expectation-law-integral}{Expectation identities and integration over the law}
\par\nobreak
Expectation inherits the preceding integral inequalities and convergence
theorems on replacing \(\mu\) by \(\mathbb P\). In particular
\(\mathbb E[\mathbf1_A]=\mathbb P(A)\), nonnegative expectations
are additive and scale by finite nonnegative constants, and integrable
expectations are linear and preserve a.e. order. If \(X\) is integrable,
\(\mathbb E[aX+b]=a\mathbb E[X]+b\) for real \(a,b\).
For nonnegative \(X\) with infinite expectation, this affine expression
still has an expectation: it is \(\infty\) for \(a>0\),
\(-\infty\) for \(a<0\), and \(b\) for \(a=0\).

Let \(X\) be a random element with law \(\mathbb P_X\) on
\((S,\mathcal S)\), and let \(h\) be measurable and either
nonnegative or real valued. Then
\[
 \mathbb E[h(X)]=\int_S h\,d\mathbb P_X
\]
whenever \(h\geq0\) or \(\mathbb E[|h(X)|]<\infty\).
Moreover \(\mathbb E[|h(X)|]=\int_S|h|\,d\mathbb P_X\), so
integrability can be checked on either space. On a countable atomic
space with masses \(p(\omega)\), the nonnegative or integrable
expectation is \(\sum_\omega X(\omega)p(\omega)\); the signed
integrable sum is absolutely convergent.
\end{hbtheorem}

\begin{hbtheorem}{thm:p15-expectation-bounds}{proved=annex:p15.expectation-bounds}{Probability bounds and variance}
\par\nobreak
Let \(X\) be a real random variable, \(A\) a nonempty Borel set,
and \(\varphi:\mathbb R\to[0,\infty)\) Borel measurable. For
\(c_A=\inf_{y\in A}\varphi(y)\),
\[
 c_A\mathbb P(X\in A)\leq
 \mathbb E[\varphi(X);X\in A]\leq\mathbb E[\varphi(X)].
\]
For \(A=\varnothing\), the event and its restricted expectation are
zero and no infimum needs to be assigned. In particular, for nonnegative
\(Y\) and \(r>0\), Markov's bound is
\(\mathbb P(Y\geq r)\leq\mathbb E[Y]/r\).

If \(\mathbb E[X^2]<\infty\), then \(X\) is integrable and
\[
 \begin{aligned}
 0\leq\operatorname{Var}(X)&=\mathbb E[X^2]-(\mathbb E[X])^2,\\
 \operatorname{Var}(aX+b)&=a^2\operatorname{Var}(X),\\
 \mathbb P(|X-\mathbb E[X]|\geq r)&\leq
                         \operatorname{Var}(X)/r^2\qquad(r>0).
 \end{aligned}
\]
Here \(a,b\) are any real constants. The last inequality is the
usual variance form of Chebyshev's bound. Durrett calls the more general
nonnegative-function bound above Chebyshev's inequality as well.
\end{hbtheorem}

\begin{hbexample}{def:p15-computed-expectations}{An integral computed without a density, and an undefined signed mean}
For the constructed uniform variable \(U(u)=u\), partition \((0,1)\)
by \((k/m,(k+1)/m]\cap(0,1)\), \(0\leq k<m\). Each cell has
measure \(1/m\). Bounding \(U\) by its two endpoints gives
\[
 \frac{m-1}{2m}\leq\mathbb E[U]\leq\frac{m+1}{2m}.
\]
Bounding \(U^2\) by the squared endpoints gives lower and upper sums
\[
 \frac{(m-1)m(2m-1)}{6m^3}
 \leq\mathbb E[U^2]\leq
 \frac{m(m+1)(2m+1)}{6m^3},
\]
both tending to \(1/3\).
The finite identities for sums of integers and squares, the two
simple-function bounds and the limits are proved in
\cref{thm:p15-w02b-p15-001-uniform-moments}.
Thus \(\mathbb E[U]=1/2\), \(\mathbb E[U^2]=1/3\), and
\(\operatorname{Var}(U)=1/12\). The calculation uses the already
constructed interval measure and simple integral bounds.

For contrast, take the countable outcomes \((s,n)\), where
\(s\in\{-1,1\}\) and \(n\geq1\), with mass \(2^{-n-1}\)
at each point. The total mass is \(\sum_{n\geq1}2^{-n}=1\).
Let \(X(s,n)=s2^n\). Each positive and each negative level contributes
\(1/2\) to its respective part expectation, so
\(\mathbb E[X^+]=\mathbb E[X^-]=\infty\). The signed mean is
undefined. Pairing the two signs and obtaining zero would be an
unauthorized cancellation of those two infinite nonnegative sums.
\end{hbexample}

\begin{hbtheorem}{thm:p15-moment-controlled-convergence}{proved=annex:p15.moment-controlled-convergence}{Convergence controlled by a stronger moment}
\par\nobreak
Let real random variables \(X_n\to X\) a.s. Let
\(g,h:\mathbb R\to\mathbb R\) be continuous and satisfy
\[
 g\geq0,\qquad g(x)\to\infty\text{ as }|x|\to\infty,\qquad
 \frac{|h(x)|}{g(x)}\to0\text{ as }|x|\to\infty.
\]
The ratio is used only sufficiently far out that \(g(x)>0\).
If \(\mathbb E[g(X_n)]\leq K<\infty\) for every \(n\), then
\(h(X_n)\) and \(h(X)\) are integrable and
\(\mathbb E[h(X_n)]\to\mathbb E[h(X)]\).
For example, \(g(x)=|x|^p\), \(p>1\), and \(h(x)=x\) give
convergence of means under a uniform \(p\)-th absolute-moment bound
and almost-sure convergence.
\end{hbtheorem}

\begin{proofboundary}{def:p15-integration-next-boundary}{Further representation and probability results}
The construction and all theorems stated here have complete annex proofs.
Product measures and Fubini are developed in \cref{sec:p15-product-measures};
the exponential, gamma and normal calculations in
\cref{sec:p15-convolution-gamma-normal}; and density representation in
\cref{sec:p15-radon-nikodym}.
\Cref{sec:p15-conditional-expectation} supplies the conditional calculus,
the \(L^p\) norm and contraction results, and the complete \(L^2\)
projection construction. The interval, step and continuous
approximations and the Riemann--Lebesgue lemma are proved in
\cref{thm:p15-criteria-real-integration}. The dyadic, truncation,
small-set and simple-function approximations are proved in
\cref{thm:p15-w02b-p15-001-approximations}; the essential bound, the
limit of \(p\)-norms, Minkowski's inequality and the normed space
\(L^p(\mu)\) on a general measure space in
\cref{thm:p15-w02b-p15-001-lp-results}; and the Poisson and geometric
moment calculations in
\cref{thm:p15-w02b-p15-001-poisson-moments,thm:p15-w02b-p15-001-geometric-moments}.
\end{proofboundary}
